\documentclass[fontset=fandol]{ctexart}

\usepackage[a4paper,margin=18mm]{geometry}
\usepackage{xcolor}

\definecolor{ctexblue}{HTML}{2563EB}
\definecolor{ctexgreen}{HTML}{15803D}
\definecolor{ctexorange}{HTML}{C2410C}
\definecolor{ctexgray}{HTML}{475569}

\newcommand\casebox[2]{%
  \par\medskip
  \noindent
  \fcolorbox{#1}{#1!4}{%
    \begin{minipage}{0.93\linewidth}
      #2
    \end{minipage}}%
  \par\medskip
}

\begin{document}

\pagestyle{empty}

{\LARGE\bfseries \textcolor{ctexblue}{CTeX \texttt{autoindent}：零缩进例外}\par}
\smallskip
{\color{ctexgray}
三组示例都启用 \verb|autoindent=3| 并切换到三号字。
字号变化前的 \verb|\parindent| 是否为零，决定了自动调整是否执行。\par}

\ctexset{autoindent=3}

\casebox{ctexgreen}{%
  {\large\bfseries\textcolor{ctexgreen}{1. 当前缩进非零：跟随字号更新}\par}
  \smallskip
  \begingroup
    \parindent=2em
    \zihao{3}
    {\small\ttfamily
      变化后 \string\parindent=\the\parindent，
      3\string\ccwd=\the\dimexpr3\ccwd\relax\par}
    \par
    这是一段采用三号字的示例文字。它的首行按新的汉字宽度缩进三个字，
    因为字号变化前的段首缩进并不是零。
  \endgroup
}

\casebox{ctexorange}{%
  {\large\bfseries\textcolor{ctexorange}{2. 用户明确设为零：继续保持零缩进}\par}
  \smallskip
  \begingroup
    \parindent=0pt
    \zihao{3}
    {\small\ttfamily 变化后 \string\parindent=\the\parindent\par}
    \par
    这是一段采用三号字的示例文字。即使 \texttt{autoindent} 已启用，
    明确的零缩进也不会因字号变化而被重新打开。
  \endgroup
}

\casebox{ctexblue}{%
  {\large\bfseries\textcolor{ctexblue}{3. minipage 等结构：不会意外引入缩进}\par}
  \smallskip
  \begingroup
    \zihao{3}
    {\small\ttfamily minipage 内变化后 \string\parindent=\the\parindent\par}
    \par
    这个彩色框本身就是一个 \texttt{minipage}；该结构把段首缩进设为零。
    字号变化后仍保持零，因此不会破坏结构原有的版式。
  \endgroup
}

\bigskip
\noindent
\textcolor{ctexgray}{结论：\texttt{autoindent} 让已有的非零缩进跟随字号变化，
而不是无条件给 \texttt{\string\parindent} 赋值。}

\end{document}
